% Matousek Partition Tree -- Short Report
% Compile: xelatex partition_tree_short_report.tex (twice for refs)

\documentclass[10pt,a4paper]{article}

\usepackage{fontspec}
\usepackage{unicode-math}
\setmainfont{Helvetica Neue}
\setmathfont{STIX Two Math}
\setmonofont[Scale=0.85]{Menlo}

\usepackage[margin=1.7cm]{geometry}
\usepackage{booktabs}
\usepackage[dvipsnames]{xcolor}
\usepackage{xurl}
\usepackage{hyperref}
\usepackage{microtype}
\usepackage{enumitem}
\usepackage{float}

\hypersetup{
  colorlinks=true,
  linkcolor=blue!60!black,
  citecolor=blue!60!black,
  urlcolor=blue!60!black,
}

\setlist[itemize]{nosep,leftmargin=*}
\setlist[enumerate]{nosep,leftmargin=*}
\setlength{\parindent}{0pt}
\setlength{\parskip}{4pt}

\title{Mato\v{u}\v{s}ek Partition Tree: Short Technical Report}
\author{}
\date{}

\begin{document}
\maketitle
\vspace{-1.2em}

\section*{Goal}

This project implements one constructive step of Mato\v{u}\v{s}ek's
2D Partition Theorem for halfplane range counting. The theorem says
that $n$ planar points can be split into about $r$ triangular groups
so that every query line crosses only
\[
  O(\sqrt r)
\]
groups. A query can count non-crossed groups wholesale and recurse
only into crossed groups.

The theorem is not new here, and neither is its cutting machinery,
which is due to Chazelle and Friedman. The contribution is
implementation and measurement: build the construction, check the
proof invariants at runtime, and measure the constants hidden in the
asymptotic bound.

\section*{Principle}

The construction is easiest to understand as a way to avoid checking
all possible query lines.

There are infinitely many lines, so the algorithm first builds a
finite test set $Q$. It then constructs the point groups while keeping
the crossing count small for every line in $Q$. A separate lemma, the
Test Set Lemma, transfers this finite guarantee to all query lines.

So the proof has two levels:

\begin{enumerate}
  \item Control the finite test set:
        \[
          K_Q=O(\sqrt r).
        \]
  \item Transfer that bound to every line:
        \[
          \operatorname{cr}_\Pi(h)\le 3K_Q+\sqrt r.
        \]
\end{enumerate}

\section*{Algorithm Summary}

The construction has three main ingredients.

\textbf{Duality.} Each point $p=(a,b)$ is mapped to the line
$p^*:y=ax-b$. This turns reasoning about query lines into reasoning
about cells in the dual plane.

\textbf{Cuttings.} A cutting divides the dual plane into cells that
are crossed by only a controlled number of lines. A first cutting
builds a finite test set $Q$ with
\[
  |Q|\le r.
\]
This implementation uses cutting scale $t=\beta\sqrt r$ with
$\beta=0.25$, and checks the size condition directly.

\textbf{Reweighting.} Each test line starts with weight $1$. Whenever
a chosen group triangle crosses a test line, that line's weight
doubles. Future weighted cuttings avoid high-weight lines. This keeps
the worst test-line crossing count
\[
  K_Q=\max_{q\in Q}\kappa(q)
\]
near a constant multiple of $\sqrt r$.

\section*{Correctness Checks}

The code raises an error instead of silently continuing when a proof
condition fails. The main runtime checks are:

\begin{itemize}
  \item every cutting cell satisfies the crossing-weight bound;
  \item the number of cutting faces stays within the pigeonhole budget;
  \item the test set satisfies $|Q|\le r$;
  \item groups are disjoint and contain the required points;
  \item exact query answers match brute force tests.
\end{itemize}

All geometric predicates use exact rational arithmetic through
\texttt{fractions.Fraction}. These checks are not a formal proof
certificate, but they make the implementation auditable.

\section*{Experiment}

The experiment uses $n=1200$ uniformly random points with seed 7 and
applies one partition step, not the full recursive tree:
\[
  \texttt{python3 benchmarks/measure\_crossings.py 1200 7}.
\]

\begin{table}[H]
\centering
\caption{Crossing measurements for one verified partition step.}
\begin{tabular}{@{}rrrrrrl@{}}
\toprule
$r$ & groups & $|Q|$ & $K_Q$ & $K_Q/\sqrt r$ &
  random max & $3K_Q+\sqrt r$ \\
\midrule
25  & 25 & 17 & 20 & 4.00 & 25 & 65 (vacuous) \\
36  & 36 & 31 & 25 & 4.17 & 34 & 81 (vacuous) \\
64  & 66 & 55 & 39 & 4.88 & 48 & 125 (vacuous) \\
\bottomrule
\end{tabular}
\end{table}

\vspace{-0.6em}

\section*{Takeaway}

The mechanism works: in these tests, $K_Q$ grows like $\sqrt r$, with
measured constant about $4$--$5$.

The constants are large. The Test Set Lemma transfers the test-line
bound to all lines using
\[
  \operatorname{cr}_\Pi(h)\le 3K_Q+\sqrt r.
\]
With the measured constants, this bound is still larger than the
total group count until roughly $r=170$. At that scale, the exact
rational construction is already slow.

So the project demonstrates the theorem's mechanism clearly, but also
shows why this theory-optimal construction is hard to turn into a
practical data structure at moderate input sizes.

\section*{Repository}

Code, tests, benchmarks, the full technical report, and the longer
derivation are available at:
\url{https://github.com/wangyi1010/matousek-partition-tree}.

\end{document}
